% versions-bank.tex -- one problem that reports which version it was typeset
% for, drawn by a versioned exam (versions-test.tex), an exam with no \versions
% (versions-none-test.tex) and a quiz (versions-none-quiz.tex).
%
% \onver{label}{content} is the course-side selector \IfExamVersioned exists
% for: it prints content on the copy whose version is label, and a document with
% no versions prints the first label's branch. The comparison with
% \theExamVersion sits inside the versioned branch because that macro is A in an
% exam with no \versions and undefined in a quiz, so a selector that compares
% without asking prints nothing in the first and raises "Undefined control
% sequence" in the second.
\ExplSyntaxOn
\NewDocumentCommand{\onver}{ m m }{
	\IfExamVersioned
		{ \str_if_eq:eeT { #1 } { \theExamVersion } { #2 } }
		{ \str_if_eq:nnT { #1 } { 1301 } { #2 } }
}
\ExplSyntaxOff

% VQFLAG reads the query through \setvar, whose argument is expanded inside
% \directlua: it prints 1 or 0 only while \IfExamVersioned stays expandable.
% VQLABEL prints \theExamVersion as the class left it.
\begin{problem}{vquery}[topic=fr]
	\setvar{vqflag}{\IfExamVersioned{1}{0}}
	VQSTEM \IfExamVersioned{VQVERSIONED}{VQUNVERSIONED}
	\onver{1301}{VQSEL1301}\onver{1302}{VQSEL1302}\onver{1303}{VQSEL1303}\onver{1304}{VQSEL1304}\onver{1305}{VQSEL1305}\onver{1306}{VQSEL1306}
	VQFLAG\get{vqflag}
	VQLABEL[\ifdefined\theExamVersion\theExamVersion\else none\fi]
\end{problem}
